\section{Loop and Tadpole}
\label{sec:res_scp}
% ======================================================================

Beyond the dynamic three-phonon bubble $\Sigma_B(\omega)$ (the only term
with an imaginary part, hence the linewidth), the lowest-order diagrammatic expansion
adds two static, real, Hermitian self-energies: the quartic loop
$\Sigma_L=\tfrac12\,\Phi_4\!:\!\langle uu\rangle$ (needs FC4) and the cubic tadpole
$\Sigma_T=\Phi_3\!:\!\langle u\rangle$ with
$\langle u\rangle=-\Phi_\mathrm{eff}^{-1}\,\tfrac12\Phi_3\!:\!\langle
uu\rangle$ (internal-coordinate relaxation, needs FC3). We computed all
three contributions across temperature for the soft-mode
\texttt{d5a} nanowire and the stiff (3,3) carbon nanotube
(\cref{fig:res_scp_mag}). The CNT FC4 was obtained from the existing
rattled set with a short ($2.0$~\AA) quartic cutoff,
the \texttt{d5a} FC4 from the \texttt{rfe} refit of \cref{sec:res_fc}.

For the CNT, the dynamic bubble dominates,
$\max|\mathrm{Re}\,\Sigma_B|:\|\Sigma_L\|:\|\Sigma_T\|\approx 8:1:0.5$
at $300$~K, all growing with $T$ as $\langle uu\rangle$ grows. For the
\texttt{d5a} wire, the ordering inverts: the bubble is weak
($\max|\mathrm{Re}\,\Sigma_B|\!\approx\!670$~THz$^2$ at $300$~K), while the
loop is the largest correction
($\|\Sigma_L\|\!\approx\!1045$~THz$^2$). In both structures, the loop
exceeds the tadpole at every temperature.

\begin{figure}[htb]
  \centering
  \includegraphics[width=0.95\linewidth]{transport_sweeps/static_se_study.pdf}
  \caption{Static-correction magnitudes (top) and heat-flow conservation
  per correction (bottom) versus temperature, for the CNT (left)
  and the \texttt{d5a} nanowire (right)}
  \label{fig:res_scp_mag}
\end{figure}

For the CNT, the loop is a $\sim10\%$ renormalisation that does not improve the
($10^{-4}$) conservation, so it is unnecessary. For \texttt{d5a}, it is the
dominant static term. It makes the high-temperature SCBA converge, and it improves the
heat-flow conservation (halving it from $2.98\%$ to $1.34\%$ at $30$~K,
and reaching $0.45\%$ at $100$~K with the loop$+$tadpole). The loop is the
physical cure for a soft mode: stiffening it lifts the propagator out of
the near-singular region.

Run in isolation, the tadpole diverges: for \texttt{d5a} the
tadpole-only self-energy reaches $\|\Sigma_T\|=1298$~THz$^2$ at $300$~K and drives the lowest mode imaginary ($\omega^2=-2.10$), with the SCBA failing and the heat-flow
conservation increasing. The optical tadpole
$T_O$ is the gradient that relaxes the internal coordinates and
vanishes at the relaxed structure~\cite{paulatto_anharmonic_2013}.
With coupled self-consistency, the quartic loop
stiffens the modes, which lowers $\langle uu\rangle$, which shrinks
$\langle u\rangle$ and hence $\Sigma_T$. Running loop and tadpole together
(as SSCHA does~\cite{tadano_first-principles_2018}) the same $300$~K point gives
$\|\Sigma_T\|=284$~THz$^2$ ($4.6\times$ smaller) and a stable
lowest mode ($0.765$~THz). On bulk silicon, $\|\Sigma_T\|$ stays at $1$--$7$~THz$^2$
(a few percent of the optical $\omega^2$) across $100$--$600$~K and never
destabilises.

\begin{figure}[htb]
  \centering
  \includegraphics[width=0.95\linewidth]{transport_sweeps/static_se_tadpole_breakdown.pdf}
  \caption{Tadpole magnitudes versus temperature for the CNT and
  \texttt{d5a}. Circled points mark where the tadpole-only self-energy
  drives an imaginary renormalised mode and the SCBA does not converge.}
  \label{fig:res_scp_tadpole}
\end{figure}

% ======================================================================
